\section{Selecting a tree from relaxation information}
\label{sec:branching}

A small certificate does not imply that a branching rule will find a small
tree. This section fixes the oracle and compares the tree selected by a
rule with the smallest permitted certificate. The comparison counts bound
evaluations. It does not count the work needed to compute a bound, discover
an optimal incumbent, query a function germ, or compute an optimal cut.

Here $F=D=\prod_i[L_i,U_i]$, $f$ is continuous, the incumbent is fixed at
$f_* =\min_Df$, and $\epsilon>0$. For a fixed $\alpha>0$, the oracle is
\begin{equation}
 f_C(y)=f(y)-\alpha q_C(y),\qquad L(C)=\min_Cf_C,
 \qquad q_C(y)=\sum_i(y_i-l_i)(u_i-y_i).
 \label{branching:oracle}
\end{equation}
Put $M_\epsilon(y)=f(y)-f_*+\epsilon$. A box is valid exactly when
$M_\epsilon(y)\ge\alpha q_C(y)$ everywhere on $C$. At each invalid
node a binary rule makes a strictly interior coordinate cut and processes
both children. Thus $T=2N-1$ for a finished tree with $N$ leaves.
One exact bound evaluation is charged per node. Validity is hereditary,
because $q_{C'}\le q_C$ on $C'\subseteq C$. Uniformly small boxes are
valid, since $M_\epsilon\ge\epsilon$.

The minimum rectangular-partition size is $N_{\rm rect}$, and the minimum
guillotine-tree leaf count is $N_{\rm tree}$. In one dimension they agree,
and $T_{\rm opt}=2N_{\rm rect}-1$. In higher dimensions a partition may
require refinement to be realized by a tree. No convexity of $f$ is
assumed below. Even convexity of the underestimator in
\eqref{branching:oracle} is unnecessary for the one-dimensional theorem.
If $f+\alpha\|y\|^2$ is convex, the oracle is a convex minimization
oracle; otherwise exact minimization may itself be expensive.

\begin{theorem}[One-dimensional minimizer branching]
\label{branching:four-competitive}
On an interval, split each invalid node at any exact minimizer of
$f_C$. For every valid partition with $N\ge2$ intervals, the rule has
at most $4N-5$ splits and $8N-9$ nodes. If $N=1$, the root is valid
and the count is one. Consequently
\[
 T\le4T_{\rm opt}-5\quad\text{when the root is invalid},
 \qquad T\le4T_{\rm opt}\quad\text{on every instance}.
\]
The rule terminates, and its bound permits arbitrary choices among tied
minimizers, including choices using previous search history.
\end{theorem}

\begin{proof}
Use the shifted relaxation
$\phi_C=M_\epsilon-\alpha q_C$. On an invalid interval its minimum is
negative, while endpoint values are positive. Every selected minimizer
is therefore strictly interior. Split points cannot repeat: an ancestor
split is a descendant endpoint, and nonnested nodes have disjoint interiors.

Fix a certificate interval $J=[s,s']$ of length $\lambda$. An invalid
node whose split lies inside $J$ cannot be contained in $J$.
Suppose nested split nodes $C_1\supsetneq C_2$ both cross $s$, and their
splits $y_1,y_2$ lie in $(s,s')$. The child containing $s$ has upper
endpoint $y_1$, so $y_2<y_1$. Write
$d=s-l_{C_1}>0$, $e=u_{C_1}-s>0$, and $t_k=y_k-s$.
Optimality at $y_1$, invalidity at $y_2$, and validity of $J$ give
\begin{align*}
 M_\epsilon(y_1)&\le M_\epsilon(y_2)
 +\alpha[(t_1+d)(e-t_1)-(t_2+d)(e-t_2)],\\
 M_\epsilon(y_2)&<\alpha(t_2+d)(t_1-t_2),\\
 M_\epsilon(y_1)&\ge\alpha t_1(\lambda-t_1).
\end{align*}
For the second inequality, $C_2$ is contained in $[l_{C_1},y_1]$.
Cancellation yields
$t_1(\lambda-t_1)<(t_1-t_2)(e-t_1)<t_1(e-t_1)$.
Thus $\lambda<e$: the larger node also crosses $s'$.
Reflection gives the analogous right-endpoint statement.

Nodes crossing a given endpoint are nested. The crossing argument allows
at most one split node crossing only the left endpoint of $J$, and at most
one crossing only its right endpoint. At most one crosses both, since its
split separates the two endpoints. Thus an interior certificate interval
contains at most three split points. Each end interval contains at most
one. The $N-1$ interior certificate breakpoints can each be used once.
The total is at most $(N-1)+2+3(N-2)=4N-5$.
The argument applies to every finite partial tree. An infinite run would
have a finite partial tree exceeding this bound. Continuity gives a
minimizer at every node, and a uniform valid partition supplies a finite
comparison certificate. This proves termination and the node count.
\end{proof}

\begin{proposition}[Two-piece certificates and the known ratio range]
\label{branching:ratio-range}
If $N_{\rm rect}=2$, minimizer branching uses at most five nodes, and
this count is attained. Over all continuous one-dimensional instances,
its worst-case competitive ratio lies in $[11/5,4]$.
\end{proposition}

\Cref{app:branching-small-certificates} proves the five-node refinement
and gives a fully specified eleven-node instance with a three-piece
optimum. The strict instancewise inequality in
\cref{branching:four-competitive} does not imply a supremum strictly
below four. The exact worst-case ratio remains undetermined.

We next specify what a rule knows. A \emph{local relaxation-information
rule} sees its current box, $\alpha$, $\epsilon$, incumbent, exact bound,
a selected relaxation minimizer, and the germ of $f$ at that minimizer.
A germ gives equality on some open neighborhood; it does not allow
arbitrary queries over the box. The rule carries no objective-dependent
information from other nodes. For the following root obstruction this
last restriction is immaterial, but it matters in repeated ambiguities.

\begin{proposition}[An information obstruction on a nonanalytic class]
\label{branching:germ-obstruction}
There are two one-dimensional continuous objectives, with convex
underestimators in \eqref{branching:oracle}, having identical root data,
identical germs at the root minimizer, and identical functions near both
root endpoints. Each has a three-node optimal tree, but their optimal
first cuts differ. For every deterministic rule using these data, at
least one of the two instances has ratio at least $5/3$. For every
randomized root-cut rule, at least one has expected ratio at least $4/3$.
The same statements hold in a class of $C^\infty$ nonanalytic objectives.
\end{proposition}

The construction and its smoothing proof are in
\cref{app:branching-information}. It uses nonanalytic functions explicitly.
For real-analytic objectives on a connected neighborhood, a germ determines
the whole objective by uniqueness of analytic continuation; this
information obstruction does not apply to that class. Together with
\cref{branching:four-competitive}, the deterministic minimax ratio in
the stated nonanalytic information class belongs to $[5/3,4]$.

Full function information can select an optimal certificate without
information loss. In one dimension define $b(l)$ as the largest right
endpoint for which $[l,b(l)]$ is valid. Validity is closed in the endpoint,
so the largest endpoint exists. Starting at the left root endpoint,
take this largest interval repeatedly. If $s_j$ are endpoints of any
optimal partition and $x_j$ are greedy endpoints, induction gives
$x_j\ge s_j$: either $x_{j-1}\ge s_j$, or
$[x_{j-1},s_j]$ is a valid subinterval of the $j$th comparison piece.
Thus the greedy partition has the optimal size. Sufficiently short
intervals are uniformly valid, so the construction terminates.
In higher dimensions a full function oracle can choose the first cut
of a minimum guillotine certificate, and repeat on its children. This
statement counts the selected tree, not the cost of constructing the
optimal certificate.

\subsection{Safeguards and their unavoidable cost}

A $\theta$-safe cut of an interval of width $w$ lies in
$[l+\theta w,u-\theta w]$, where $0<\theta<1/2$.
Every child then has width at most $(1-\theta)w$. Along an arbitrary
path its width halves within
$\lceil\log2/[-\log(1-\theta)]\rceil$ splits.
This protects against relaxation points at or close to a bound.
Exact minimizers at invalid nodes in the preceding one-dimensional model
are already interior, but can be arbitrarily close to an endpoint.

For $f_a(y)=2\alpha|y-a|$ on $[0,1]$, the exact minimizer on an interval
straddling $a$ is $a$, and the interval is invalid exactly when
$\alpha(a-l)(u-a)>\epsilon$. Every interval on one side of $a$ is valid.
If the root is invalid, cutting at $a$ gives $T_{\rm opt}=3$.

\begin{proposition}[Persistent clipping can lose logarithmically]
\label{branching:persistent-clamp}
For fixed $0\le\lambda\le1$ and $0\le\theta<1/2$, use the relative cut
\[
 \sigma(p)=\operatorname{clip}\left(\lambda p+(1-\lambda)/2,
                                      \theta,1-\theta\right),
\]
where $p$ is the relaxation minimizer's relative position. Unless
$(\lambda,\theta)=(1,0)$, some fixed kink $a\in(0,1/2)$ has
$T=\Omega(\log(1/\epsilon))$ although $T_{\rm opt}=3$.
\end{proposition}

\begin{proof}
The continuous function $\sigma(p)-p/(1-p)$ is positive at zero and
negative at $1/2$, so it vanishes at a point $p\in(0,1/2)$. Put
$a=p$ and $\kappa=p/(1-p)\in(0,1)$. Symmetry of $\sigma$ gives a
kink-containing chain with relative positions alternating between
$p$ and $1-p$, and width $\kappa^k$ at depth $k$. All off-chain
children are valid. The chain node is invalid when
$\alpha p(1-p)\kappa^{2k}>\epsilon$. Thus, for
$\epsilon<\alpha p(1-p)$, it uses at least
\[
 2\left\lceil\frac{\log(\alpha p(1-p)/\epsilon)}
                         {2\log(1/\kappa)}\right\rceil+1
\]
bound evaluations. The pure minimizer rule instead ends in three nodes.
\end{proof}

\begin{proposition}[The price of safety and a recentring rule]
\label{branching:recentring}
If $a<\theta$, put
$J=\lceil\log(\theta/a)/\log(1/\theta)\rceil$.
Every $\theta$-safe rule, with any information and any randomization,
has $T\ge2J+1$ on $f_a$ whenever
$\epsilon<\alpha a^2(1-\theta)/\theta$.
Consequently no fixed safety margin gives a uniformly bounded competitive
ratio over kink locations as $\epsilon\downarrow0$.

For $0<\theta\le1/3$, the following recentring rule nevertheless has a
bounded count as $\epsilon\downarrow0$ for each fixed kink. Cut at the
minimizer $p$ when it is safe; below the safe region cut at
$l+\max\{\theta w,2(p-l)\}$; above it use the reflected formula.
If $d_0=\min\{a,1-a\}$ and
$J=\min\{j\ge0:d_0\ge\theta^{j+1}\}$, it uses
$T\le2J+3$. Its zero-tolerance chain hits the kink in exactly $J+1$
splits, the minimum among safe chains that hit it.
\end{proposition}

\Cref{app:branching-safety} proves these claims. The worst-case location
in the first assertion may depend on $\epsilon$; the persistent-clamp
counterexample instead uses one fixed location. Recentring gives an
explicit tradeoff between width contraction and kink hitting. Its hitting
optimality does not assert optimality among all finite-tolerance
certificates, since a positive tolerance can prune before hitting the kink.
All counts here are bound evaluations for the fixed oracle. They are not
runtime comparisons between numerical minimization methods.

\subsection{Coordinate choice and a positive restricted extension}

For a separable objective write
$f-f_* =\sum_i m_i(x_i)$, with $\min_{D_i}m_i=0$.
Let $N_i(b)$ be the least number of intervals partitioning $D_i$ with
$\min_J(m_i-\alpha q_J)\ge-b$. Slicing a rectangular certificate at
minimizers of all other coordinates, and taking products of one-dimensional
certificates, gives
\begin{equation}
 \max_iN_i(\epsilon)\le N_{\rm rect}(\epsilon)
 \le N_{\rm tree}(\epsilon)
 \le\inf_{\epsilon_i>0,\,\sum_i\epsilon_i=\epsilon}
                                      \prod_iN_i(\epsilon_i).
 \label{branching:separable-brackets}
\end{equation}
The slice uses the whole tolerance; the product allocates it across axes.
These inequalities alone do not settle coordinate selection.

\begin{proposition}[Repeated coordinate ambiguity]
\label{branching:dimension-obstruction}
For every $n\ge2$, there are $n$ nonanalytic separable objectives on
$[0,1]^n$, at the common tolerance $\epsilon=1/100$ and $\alpha=1$,
each with $T_{\rm opt}=3$. For every deterministic node-local
relaxation-information rule, with coordinate-wise minimizer selection,
at least one of these instances satisfies
\[
 T\ge2\left\lceil\frac{2^{n+1}-n-3}{n-1}\right\rceil+1.
\]
For every distribution of such local rules, at least one instance has
\[
 \E T\ge\frac{2(2^{n+1}-n-2)}n+1.
\]
The rules may also see the objective germ at every corner of the current
box. They may not carry objective-dependent information from other nodes
or use objective information outside these germs to resolve minimizer ties.
\end{proposition}

\Cref{app:branching-dimension} supplies the family and the counting proof.
The result forces exponential dependence on dimension in any competitive
constant for this information class. It does not show accuracy-dependent
loss at a fixed dimension, or apply to analytic germs or rules that learn
from siblings.

For example, on $f(x,z)=x^2$ in $[0,1]^2$ with $\alpha=1$, a rule that
splits every coordinate with an interior relaxation minimizer has
$\Omega(\epsilon^{-1/2}\log(1/\epsilon))$ processed boxes, whereas a
guillotine certificate has $O(\epsilon^{-1/2})$ leaves.
\Cref{app:branching-multi} proves this separation. Selecting the correct
point in each coordinate does not justify splitting all coordinates.

The $\omega$ rule instead selects a coordinate maximizing
$q_{J_i}(y_i)$, where $y_i$ is its selected exact relaxation minimizer,
and cuts there. General constant competitiveness of this rule, even
for arbitrary separable objectives in a fixed dimension at least two,
remains open. The following extension isolates a class in which the
phase accounting can be completed.

\begin{theorem}[One arbitrary coordinate and sharp coordinates]
\label{branching:sharp-induction}
Scale to $\alpha=1$. There is one arbitrary continuous nonnegative
coordinate function with minimum zero on its original interval $A$,
and $r$ sharp coordinate functions with
minima zero at interior kinks $c_i$. Fix a coordinate-wise minimizer
selection depending only on its function and interval. A sharp coordinate
means that on every interval straddling $c_i$, the selected minimizer is
$c_i$ with value $-q_J(c_i)$, and on either side the selected minimizer
has zero gap and nonnegative relaxation value. Use the fixed optimal
incumbent and exact separable oracle above.
Explicitly, $f-f_*=m_0(x)+\sum_{i=1}^rm_i(z_i)$, and $N_A(b)$
is the least number of intervals partitioning $A$ for which
$\min_J(m_0-q_J)\ge-b$.

Set $C_0=4$ and, for $r\ge1$, define
\[
 c_r=r[2r(r+1)+1],\quad A_r=4(1+c_r),\quad
 C_r=(A_r+1)+2C_{r-1}(A_r+2).
\]
Then $\omega$ terminates with at most $C_rN_A(\epsilon)$ leaves and
\[
 T_\omega\le2C_rN_{\rm rect}(\epsilon)-1.
\]
Coordinate-choice ties may be resolved arbitrarily. Every fixed
coordinate-wise minimizer selection is permitted if every sharp-coordinate
minimizer satisfies the stated properties; otherwise fix a selection
satisfying them.
\end{theorem}

The complete induction and phase lemma are in
\cref{app:branching-sharp}. The constants depend only on $r$ and are
conservative. For an explicit sharp coordinate on $[L_i,U_i]$, take
$m_i(t)=\sigma_i|t-c_i|-(t-c_i)^2$, with
$\sigma_i\ge2(U_i-L_i)$ and $c_i\in(L_i,U_i)$.
The left and right slopes of $m_i-q_{[l,u]}$ are
$-\sigma_i-(l+u-2c_i)$ and $\sigma_i-(l+u-2c_i)$.
They are strictly negative and positive respectively, so the kink or the
nearest endpoint is the unique minimizer and the sharpness conditions hold.
The arbitrary coordinate may be quadratic or nonconvex. The theorem does
not cover two
arbitrary coordinates or justify a global sum over phases for general
separable objectives.
